\section{Pricing Schedules}
\label{sec:pricing}

% VALUE: the problem in words: a schedule makes three essential decisions (the split; the order in which the CP interleaves each command's segments, each operator's kept in order; the buffer assignment); operators hand data over through DRAM, so a consumer issues after its producer's last synchronization and any host work, in a separate command; D and packing manage only the host boundary and are reported apart; the objective is the makespan
% DECISION: no NP-hardness claim until a reduction is written; with negligible issue cost and one layer's fixed per-chunk pattern the problem has the structure of a two-machine chain-reentrant shop (wang1997reentrant), and with non-negligible issue cost the CP acts as a shared server (glass2000dedicated, hall2000common) (vault ch3-formulation, Relation to classical scheduling).
\label{subsec:formulation}
A schedule of a compiled design makes three essential decisions, namely the split, which sets which operators can run at once (Section~\ref{subsec:prefill}), the order in which the CP interleaves each command's segments, keeping every operator's segments in their own order, and the slot that each use of a buffer takes, since two uses of one slot run in turn even when neither consumes the other's data.
Because operators hand data over through DRAM with no lock between them, a consumer issues only after its producer's final synchronization and any host work between the two, in a separate command.
The number $D$ of commands in flight and the packing of commands manage only the boundary with the host, so we report their effect apart, and the objective is the makespan, which for prefill is the TTFT less the work before the first decoder layer and after the last.

\subsection{Bounding a Removal}
\label{subsec:removal}

% VALUE: one removal is worth at most the overlap it blocks on typed resources, given every other edge; zero off the longest path; in a loop the maximum cycle mean governs; Amdahl's law for an ordering, from the order and measured durations, not a trace
Removing a false dependence $d$ speeds a phase up by at most $T / (T - O(d))$, where $T$ is the phase's time and $O(d)$ is the time during which $d$ keeps work on disjoint resources from overlapping, given every other edge of the happens-before graph.
The resources are typed, namely the column groups, the CP with its cost at each command boundary, the DRAM read rate up to its cap, and the channels that one group borrows from another's columns, and since none overlaps its own work, $O(d)$ is zero when $d$ lies off the longest path.
In a loop, the maximum cycle mean of the repeating graph determines the time per iteration instead~\citep{ghamarian2006throughput}.
The bound is thus Amdahl's law applied to an ordering, computed from the program and the probes without a timing trace (Section~\ref{sec:related}).

% VALUE: for a set of removals the textbook critical-path condition applies (cited, one sentence, never a finding); we use it as a test that discards sets that cannot pay; on this machine it is necessary, not sufficient (the next longest path; shared capacity); the decode instance is in the evaluation
A set of removals gains only if it cuts every path or cycle that attains the maximum, which is the textbook condition of interaction-cost analysis~\citep{fields2003interaction}, and then only down to the next longest path and to what the shared DRAM read cap leaves.
We use it as a test on the recovered graph with measured costs, which discards every candidate set that cannot pay.
% SOURCE: the acyclic condition is the minimum cut of a critical subnetwork in project crashing (Kelley 1961; Fulkerson 1961), the cyclic one the iteration bound and maximum cycle mean (Ghamarian 2006; Stuijk 2008), and the set view is interaction cost (Fields 2003); all are cited in Section 8, never claimed.
% SOURCE: the test is the in-tree set test, acyclic half (longest path, cut, next longest; mlir-aie 7240529eb562, lit 19/19) and loop half (periodic graph, Howard's algorithm, critical subgraph; mlir-aie 54a8b8e192c, lit 10/10) (vault paper-cosched-decisions-log, P1 and P3).
% DECISION: the critical-path condition is textbook and stated once, here, with its citation; it is a tool, never a finding, and the decode partial-rename instance and the 8B at the cap live in Section 7 (section contract; outline v10.1).
% DECISION: the condition is stated as necessary only: a longest-path or cycle criterion is valid on an uncapacitated graph, and this machine has a shared DRAM cap, a shared command processor and a per-command boundary cost; the capacity half is the measured regime of this section, and the 8B at cap and the loss at 512-token appends are the two negative instances inside the paper.
% DECISION: the bound's sentences follow outline v9 section 3's rows (P4 withdraws them, P4u marks them untested); the critical-cycle sentences follow the W1-W5 table.

% VALUE: the bound is prescriptive (build a removal only where its bound exceeds its cost); since it never falls below one it cannot predict a loss, which the event model prices instead (the short prompts, in the evaluation)
The bound cannot predict a loss, such as the per-command costs that make reordering lose on prompts of two chunks, so the event model prices those (Section~\ref{subsec:shape-eval}).
% SOURCE: the bound and its typed resources (outline v9 section 2, C1; vault lab-bound-validation (1)).

\subsection{Event Model}
\label{subsec:pricing}

% VALUE: to choose among orders and splits without running each, an event simulation of the device, the CP, and the host prices a schedule; its form follows the platform (one in-order queue of commands, a host timeline with waits, column groups that overlap within a command) and its parameters are the operators' own timings plus a few machine constants measured once by microbenchmarks; the existing program runs through the same simulation, so both sides are priced alike; device time, host exposure, and head and tail reported apart
% COMMENT: this subsection describes the event model's rebuilt form (setup by column group, overlap bounded by bytes over the read cap, the head as a probed operator), which npu-gap is implementing in development; confirm every term against the implementation before the freeze.
To choose among orders and splits without running each, we price a schedule with an event simulation of the device, the CP, and the host, whose form follows the platform and whose parameters are measured.
Queued commands run one at a time (Section~\ref{subsec:semantics}), so command $k$ ends at
\begin{equation}
\label{eq:cmd}
e_k = \max\bigl(e_{k-1} + s_k,\; q_k + F\bigr) + b_k,
\end{equation}
where $q_k$ is the time the host queues it and $s_k$ is the transition into it, which is the setup $s_\mathrm{same}$ within one column group, and otherwise the setup $s_g$ of its group less the window $w_h$ at the end of the previous command, on group $h$, during which other columns can already start, never below zero.
Within a command, each column group's time is the sum of its operators' timings alone, with the same-group setup between consecutive operators.
When groups overlap, a group whose own reads already run at the read cap $R$ absorbs every byte its co-runners read, so the command takes
\begin{equation}
\label{eq:overlap}
b_k = \max_{g}\Bigl(t_{k,g} + [\,g \text{ at the cap}\,]\sum_{h \ne g} \frac{\beta_{k,h}}{R}\Bigr),
\end{equation}
where $t_{k,g}$ is group $g$'s time and $\beta_{k,h}$ the DRAM bytes that group $h$ reads in command $k$, counting every re-read.
Whether a group runs at the cap is measured, by comparing its bytes over $R$ with its time alone, and in the designs we evaluate the GEMMs do and attention does not.
On the host side, each command costs an issue and each operator a post-processing step, and the host waits on the oldest command once $D$ are in flight.
The existing program runs through the same simulation as its own host program, so both sides of every comparison are priced alike.
Finally, the event model adds the work before the first decoder layer and after the last, which it prices as ordinary operators, and reports the device time, the host's exposed time, and that work separately, so an error in one term cannot hide behind another.
% SOURCE: GEMM-only commands read 90-93% of the cap and attention-only about 15% during Llama-3.2-1B's 4,096-token prefill (amd-npu 71901c50 metrics/retro-2026-10-21/dram-retro.txt; vault lab-metric-fields (1) and its amendment 1); the read cap from the 8x2 bracket, 53.93 GB/s at the calibration hold (amd-npu d94266bf prefill-overlap/calib/run-2026-10-22/batch.log), re-read at every batch.
% SOURCE: the overlap form (ruling 2026-10-24): moved read bytes from the dispatched streams (aie-cosched f7eedbc, README 7e00ac9); bytes over R against time alone 0.87-1.06 for every GEMM set and about 0.15 for attention, so any threshold between 0.2 and 0.87 classifies alike; on the 30-cell co-run grid of 2026-10-22 the residual is a median 0.016 ms, the rest a fixed part of 0.030-0.041 ms carried by the transition term (npu-gap, 2026-10-24). The transition form (ruling 2026-10-24): s = s_same within a group, max(0, s_g - w_h) across groups; from the 10-22 sequences a GEMM after attention waits about 47 us and attention after a GEMM starts about 8 us early (npu-gap, device.json fit), to be pinned by microbenchmark v2. opcord/model/event_model.py simulate() (one in-order queue; host program with waits); the rebuilt form, constants schema version 2 (npu-gap, opcord dev checkout; vault decisions log "The model rebuilt from a logical base", 2026-10-24); bytes per operator instance from the instruction streams' buffer descriptors (pass-0 tools/bprime_check byte model).

% VALUE: the operators' timings and bytes depend on the LLM and the split and are measured once per configuration by running each operator alone or read from its streams; the few machine constants (dispatch floor, setup within and across groups, host issue and post, read cap) are measured once by microbenchmarks, all non-negative, none fitted to an LLM's end-to-end runs; precedent: analytical models whose parameters are machine constants measured by microbenchmarks
Only the operators' timings and bytes depend on the LLM and the split, so we measure the timings once per configuration by dispatching each GEMM alone and each attention chunk alone at several cached-context lengths, and read the bytes from the operators' streams.
Between measured points the event model interpolates, and beyond them it refuses to price.
The machine constants, namely $F$, the setups and windows of the transitions, the host's issue and post-processing costs, and $R$, are measured once by microbenchmarks on dummy data, as analytical models of communication and of memory-bound kernels measure theirs~\E{cite: Hockney; LogP; roofline}.
No constant is fitted to an LLM's end-to-end runs, so every configuration of Section~\ref{subsec:accuracy} tests the event model rather than supplying its constants.
% SOURCE: microbenchmarks (npu-gap prefill-overlap/microbench.py): dispatch floor from empty commands; setup from back-to-back sequences mixing same-group and cross-group pairs, non-negative least squares on one set and checked on another; read cap from co-runs on both sides of the crossover between the longest group and bytes over the cap; per split build (S1, S2, S4 of Llama-3.2-1B) with a Qwen3-4B transfer check; values \E{pending}.
% COMMENT: bib entries needed (read each source first): Hockney's model, LogP / LogGP, the roofline model.

% VALUE: the event model's error is stated on the difference between candidates, so its choice is right whenever two candidates differ by more than about twice that error and costs at most about that much otherwise; where the top candidates fall inside a frozen band, a short measurement at plan time settles the tie, and the unaided choice is always reported
Because a choice compares candidates, we state the event model's error on the difference between two candidates' times rather than on each time alone.
Its choice is then right whenever two candidates differ by more than about twice that error, and otherwise costs at most about that much.
Where the top candidates fall inside a band of 5\% for orders of one kind or 15\% between a composed and a shipped order, the composer runs each candidate briefly when it plans the shape and keeps the faster, and we always report the event model's unaided choice beside it (Section~\ref{subsec:decision}).
% SOURCE: measure-when-close, implemented test-first in opcord (planner.CLOSE_BAND, close_candidates, decide; test_measure_close 10/10); bands same 5%, cross 15%, tie 0.5%, one warm and three kept runs per candidate; about 19 min of device time over the whole grid, once per shape (vault decisions log, "Spine v18 drafted; amendment-1 batch and measure-when-close approved").
